% =====================================================================
% Explainable AI-Based Personalized Learning Path Recommendation System
% for University Students
%
% SPRINGER NATURE CONTRIBUTED-VOLUME (BOOK CHAPTER) VERSION
% Converted from the IEEE conference version. No number has changed.
%
% Compile with SNmult.cls from the Springer "LaTeX template for
% contributed works" package. Keep SNmult.cls in this folder.
%
% FIGURES: put fig1.png ... fig5.png next to this file. If an image is
% missing a placeholder box of the same size is drawn instead, so the
% file always compiles.
% =====================================================================

\documentclass[graybox]{SNmult}

\usepackage{type1cm}
\usepackage{makeidx}
\usepackage{graphicx}
\usepackage{multicol}
\usepackage[bottom]{footmisc}
\usepackage{amsmath,amssymb}
\usepackage{newtxtext}
\usepackage[varvw]{newtxmath}
\usepackage{url}

\makeindex

% Figure slot: uses the image when present, a placeholder box otherwise.
\newcommand{\FIGHEIGHT}{6.4cm}
\newcommand{\figslot}[3]{%
  \begin{figure}[t]\centering
  \IfFileExists{#1}{\includegraphics[width=0.92\textwidth,height=\FIGHEIGHT,keepaspectratio]{#1}}%
    {\framebox[0.92\textwidth]{\parbox[c][\FIGHEIGHT][c]{0.86\textwidth}{\centering\footnotesize\textit{#1 not found. Place the image file next to chapter.tex.}}}}%
  \caption{#2}\label{#3}\end{figure}}

\begin{document}

\title*{Explainable AI-Based Personalized Learning Path Recommendation System for University Students}
\titlerunning{Explainable Personalized Learning Path Recommendation}

\author{Abhay Pratap Singh, Akshay Guleria, Uday Chaturvedi, Gursimar Makkar and Shikha Kamal}
\authorrunning{A. P. Singh et al.}

\institute{Abhay Pratap Singh \at Department of Computer Science, Chandigarh University, Mohali, India, \email{23BCS11784@cuchd.in}
\and Akshay Guleria \at Department of Computer Science, Chandigarh University, Mohali, India, \email{23BCS10517@cuchd.in}
\and Uday Chaturvedi \at Department of Computer Science, Chandigarh University, Mohali, India, \email{23BCS10807@cuchd.in}
\and Gursimar Makkar \at Department of Computer Science, Chandigarh University, Mohali, India, \email{23BCS11455@cuchd.in}
\and Shikha Kamal \at Department of Computer Science, Chandigarh University, Mohali, India, \email{shikha.e12552@cumail.in}}

\maketitle

% --- online abstract (appears on SpringerLink) -----------------------
\abstract*{Digital learning platforms generate large volumes of student interaction data, yet
educational recommender systems offer little transparency, rarely model forgetting, and are
seldom evaluated in a way that shows which components actually contribute. This chapter
proposes an Explainable AI (XAI) based learning path recommender combining a
knowledge-graph-enhanced Deep Knowledge Tracing module, a prerequisite-constrained planner
and a priority-backtracking explanation module. On the Open University Learning Analytics
Dataset (OULAD), under five-fold cross-validation grouped by student, it attains AUC-ROC
$0.9538 \pm 0.0019$ and RMSE $0.2621 \pm 0.0026$, exceeding DKT ($p = 0.0009$) and a GNN
baseline ($p = 0.022$) after Holm correction. The central finding concerns evaluation:
removing the knowledge graph changes accuracy by only $-0.0001$ ($p = 0.77$), because
accuracy is scored at assessment positions and only 89 of 237 concepts carry assessments.
Under a cold-concept protocol that withholds assessments for a subset of concepts and scores
only those, the same graph improves accuracy by $+0.0325$ ($p = 0.0003$, Cohen's $d = 5.06$),
winning on every fold. The planner makes zero prerequisite violations across 906 decisions
evaluated against real learner behaviour, and explanations are assessed by objective
fidelity, sufficiency and stability rather than subjective rating.
\keywords{Explainable AI $\cdot$ Personalized learning path $\cdot$ Knowledge tracing $\cdot$ Knowledge graphs $\cdot$ Educational data mining $\cdot$ Ablation methodology}}

% --- printed abstract ------------------------------------------------
\abstract{Digital learning platforms generate large volumes of student interaction data, yet
educational recommender systems offer little transparency, rarely model forgetting, and are
seldom evaluated in a way that shows which components actually contribute. This chapter
proposes an Explainable AI (XAI) based learning path recommender combining a
knowledge-graph-enhanced Deep Knowledge Tracing module, a prerequisite-constrained planner
and a priority-backtracking explanation module. On the Open University Learning Analytics
Dataset (OULAD), under five-fold cross-validation grouped by student, it attains AUC-ROC
$0.9538 \pm 0.0019$ and RMSE $0.2621 \pm 0.0026$, exceeding DKT ($p = 0.0009$) and a GNN
baseline ($p = 0.022$) after Holm correction. The central finding concerns evaluation:
removing the knowledge graph changes accuracy by only $-0.0001$ ($p = 0.77$), because
accuracy is scored at assessment positions and only 89 of 237 concepts carry assessments.
Under a cold-concept protocol that withholds assessments for a subset of concepts and scores
only those, the same graph improves accuracy by $+0.0325$ ($p = 0.0003$, Cohen's $d = 5.06$),
winning on every fold. The planner makes zero prerequisite violations across 906 decisions
evaluated against real learner behaviour, and explanations are assessed by objective
fidelity, sufficiency and stability rather than subjective rating.
\keywords{Explainable AI $\cdot$ Personalized learning path $\cdot$ Knowledge tracing $\cdot$ Knowledge graphs $\cdot$ Educational data mining $\cdot$ Ablation methodology}}

% =====================================================================
\section{Introduction}
\label{sec:intro}

Higher education's embrace of digital learning has opened new possibilities for
individualized instruction, yet conventional courses still follow one path for all students
despite differences in prior knowledge and pace, causing course mis-selection and knowledge
gaps \cite{li2026survey,aher2013combination,kuzilek2017oulad}. Personalized learning path
recommendation can address this, but most systems concentrate on prediction accuracy and
offer little explanation for any recommendation \cite{corbett1994knowledge}. Designing a
system that describes a learner's evolving knowledge, sequences activities well and explains
its choices is therefore a major challenge. Knowledge tracing methods assess learning rather
than optimize progression \cite{piech2015dkt,zhang2017dkvmn}, automated sequencing remains
opaque to students and advisers \cite{ghosh2020akt,shen2021lpkt}, and the link between
knowledge graphs, sequential decision-making and temporal modeling is underexplored in
education \cite{kipf2017gcn,mnih2015dqn}.

A further limitation motivates this work. Such systems combine several components, each
justified by an ablation showing that its removal degrades accuracy. When a structural prior
exists to support concepts whose supervision is scarce, but accuracy is measured only where
supervision is abundant, the ablation tests the component in the one regime where it cannot
matter. The contributions of this work are:

\begin{enumerate}
\item A knowledge-graph-enhanced knowledge tracing model with a learned per-head temporal
decay, attaining AUC-ROC $0.9538$ and calibration error $0.0054$ on OULAD.
\item Evidence that conventional ablation cannot detect a knowledge graph's contribution,
and a cold-concept protocol that can ($+0.0325$, $d = 5.06$).
\item A prerequisite-constrained planner whose pedagogical guarantee is enforced by
construction, verified at zero violations across 906 decisions.
\item Objective explanation metrics (fidelity, sufficiency, stability) in place of
subjective rating, released with a reproducible pipeline.
\end{enumerate}

% =====================================================================
\section{Literature Review}
\label{sec:lit}

Early educational recommenders used content-based and collaborative filtering, which
suffer cold-start problems and data sparsity \cite{li2026survey,aher2013combination};
hybrid models improved robustness but forecast preferences rather than modeling knowledge
acquisition \cite{corbett1994knowledge}. Deep Knowledge Tracing (DKT) \cite{piech2015dkt}
learns the latent knowledge state directly from interaction data, unlike Bayesian Knowledge
Tracing's static parameters \cite{zhang2017dkvmn}, and attention-based successors
\cite{pandey2019sakt,ghosh2020akt} attend to any earlier interaction
\cite{vaswani2017attention}. Reinforcement learning frames path recommendation as a Markov
Decision Process and, combined with knowledge tracing, has been reported to reduce dropout
\cite{shen2021lpkt,zhao2026rlkt}. Knowledge graphs support structured, explainable
recommendation by modelling concepts and their dependencies
\cite{kipf2017gcn,zhang2023finegrained}, with Bloom's taxonomy providing a theoretical
foundation for path design through six progressive cognitive levels
\cite{khosravi2022xaied}, and graph neural networks handling the relational structure such
graphs induce. Multi-objective hierarchical frameworks have optimized dropout risk, learner
satisfaction and completion rates, adding priority-backtracking interpreters for
interpretability \cite{lundberg2017shap,ribeiro2016lime}, and hybrid neuro-symbolic and
generative approaches have begun to explore understandable explanation generation
\cite{yang2025explainable}.

\subsection{Research Gap}
\label{subsec:gap}

Four gaps remain. Explainability is added as an afterthought \cite{xie2025profile}, and
SHAP-based explanations dominate in education but are applied mainly at the global level
rather than to individual recommendations \cite{zhao2026rlkt}. Where explanations are
evaluated at all, it is usually by participant rating, which measures whether readers found
an explanation agreeable rather than whether it is faithful to the model. Many frameworks
ignore the dynamic evolution of knowledge, including forgetting, which compromises the
pedagogical integrity of suggested trajectories \cite{boujmiraz2026review}. Knowledge graphs
are rarely applied to sequential decision-making and explanation \cite{qiang2026xaiedu}.
Finally, the contribution of such structural priors is established by ablation on aggregate
accuracy, without checking whether that metric can observe the component at all.

Together these gaps limit customization, transparency and effective learning-path
optimization for university students. Existing systems rarely integrate knowledge graphs,
temporal knowledge tracing, constrained path optimization and pedagogically meaningful
explanation in one framework, and the validation of such systems reports component
contributions on metrics that may be unable to detect them, and explanation quality by
opinion rather than measurement.

% =====================================================================
\section{Methodology}
\label{sec:method}

\figslot{fig1.png}{Overall XAI-based personalized learning recommendation framework}{fig:framework}

\subsection{Overall Framework}
\label{subsec:framework}

The system estimates a learner's changing knowledge state from interaction history; the
prerequisite relationships of a knowledge graph then constrain a planner that selects the
next learning task; finally each recommendation receives a pedagogical explanation from a
priority-backtracking module \cite{zhang2023finegrained} (Fig.~\ref{fig:framework}).
The framework therefore has three components: learner knowledge estimation, sequential path
optimization and explanation production \cite{piech2015dkt}. Each is evaluated separately,
because a gain in one does not imply a gain in another: better prediction need not yield
better paths, and a faithful explanation need not be a complete one.

\begin{table}[t]
\caption{Dataset and derived concept graph}
\label{tab:dataset}
\begin{tabular}{p{7.2cm}r}
\hline\noalign{\smallskip}
\textbf{Quantity} & \textbf{Value} \\
\noalign{\smallskip}\svhline\noalign{\smallskip}
Enrolments / distinct learners        & 32{,}593 / 28{,}785 \\
Clickstream events                    & 10{,}655{,}280 \\
Assessment records                    & 173{,}912 \\
Label balance (success / not)         & 67.9\% / 32.1\% \\
Concepts / prerequisite / co-requisite edges & 237 / 724 / 302 \\
Concepts carrying assessments         & 89 of 237 \\
Week-placement agreement (Spearman)   & $\rho = 0.975$ \\
\noalign{\smallskip}\hline\noalign{\smallskip}
\end{tabular}
\end{table}

\subsection{Data Acquisition and Concept Graph}
\label{subsec:data}

Evaluation uses OULAD \cite{kuzilek2017oulad}: 32{,}593 enrolments, 10{,}655{,}280
clickstream events and 173{,}912 assessment records across seven modules. Records are
ordered chronologically, categorical attributes one-hot encoded and numerical attributes
standardized, and missing assessment scores are median-imputed only where a submission
exists. Timestamps are day-granular with no intra-day resolution, so a session is one
learner-day rather than an inactivity threshold. The processed data is split 80:10:10.

OULAD supplies no concept annotations or prerequisites, so the graph is derived. A concept
is a (module, week-of-study) pair. Week placement uses the \texttt{week\_from} field where
present and otherwise the median access day; where both exist they agree at Spearman
$\rho = 0.975$. Edges link consecutive weeks, link the four preceding weeks to each
assessment week, and are mined from behaviour, accepted only where ordering is consistent
and later success is measurably higher:
\begin{equation}
\begin{split}
\mathrm{prec}(a,b) &= \frac{\left|\{u : \tau_u(a) < \tau_u(b)\}\right|}{\left|\{u : a, b \in h_u\}\right|},\\[2pt]
\mathrm{lift}(a,b) &= \frac{P(y_b \mid a \prec b)}{P(y_b)},
\end{split}
\label{eq:mining}
\end{equation}
where $\tau_u(\cdot)$ is learner $u$'s first interaction with a concept, with support above
100 learners, $\mathrm{prec} \geq 0.75$ and $\mathrm{lift} > 1.1$. Requiring lift removes
pairs separated only by the academic calendar. The result is 237 concepts, 724 prerequisite
and 302 co-requisite edges, acyclic without intervention.

Two properties of the data matter for interpretation. OULAD records only \emph{submitted}
assessments: binarising at the pass mark gives a $95.6\%$ positive rate, with the fifth
percentile of submitted scores equal to the pass mark itself, and 150{,}013 of
323{,}925 expected submissions never occurred. Treating non-submission as failure gives a
$67.9/32.1$ balance and makes the target successful completion. Second, assessments exist
for only 89 of the 237 concepts; the remaining 148 are reached only through the graph.

\figslot{fig2.png}{Knowledge graph enhanced knowledge tracing}{fig:kt}

\subsection{Knowledge Graph-Enhanced Knowledge Tracing}
\label{subsec:kt}

A Transformer \cite{vaswani2017attention} encodes the learner's interaction sequence, while
a Graph Convolutional Network \cite{kipf2017gcn} encodes prerequisite relationships into
concept embeddings recomputed at every forward pass, so gradients reach the graph structure
(Fig.~\ref{fig:kt}). The combined representation lets the model account for both the
learner's past progress and the structural links between concepts. Forgetting is
modelled as an attention bias that decays with elapsed time,
\begin{equation}
b^{(h)}_{ij} = -\,\mathrm{softplus}(\gamma_h)\,\log(1 + \Delta_{ij}),
\label{eq:decay}
\end{equation}
where $\Delta_{ij}$ is the days between interactions $i$ and $j$ and $\gamma_h$ is learned
per head; $\gamma_h = 0$ gives the temporal ablation. Clickstream events update the hidden
state without contributing loss, and binary cross-entropy is computed only at assessment
positions, which makes a length-200 sequence model trainable when learners average 8.7
assessments. Because mastery is shown to learners, probabilities are calibrated by
temperature scaling \cite{guo2017calibration}, fitted on half of each held-out fold.

\figslot{fig3.png}{Prerequisite-constrained learning path optimization}{fig:planner}

\subsection{Prerequisite-Constrained Path Optimization}
\label{subsec:planner}

Recommendation is formulated as a Markov Decision Process whose state holds the mastery
vector, concepts already recommended, learner features and remaining budget, and the reward
combines learning progress, prerequisite satisfaction, a redundancy penalty and a step cost.
Rather than expecting a policy to learn the pedagogical constraint, it is encoded in the
admissible actions: a concept qualifies only if the learner has not mastered it, every
prerequisite is at the readiness threshold, and it belongs to an enrolled module, so an
unprepared recommendation is impossible by construction (Fig.~\ref{fig:planner}).

The evaluated policy is greedy. Each admissible concept is simulated as studied and compared
against an idle period of equal length, and the concept with the largest total improvement
in log-odds is chosen. The idle baseline is essential because studying advances the clock
and \eqref{eq:decay} then decays every earlier interaction; log-odds are needed because
engaged learners saturate near $0.98$. Thresholds are the 70th and 30th percentiles of each
learner's own mastery, since mastery is bimodal. A learned policy is left to future work.

\figslot{fig4.png}{Explainable AI module for recommendation justification}{fig:xai}

\subsection{Explainable Recommendation}
\label{subsec:xai}

The priority-backtracking interpreter traverses the chosen concept's prerequisites, scoring
each by its shortfall below readiness, discounted by graph distance (Fig.~\ref{fig:xai}).
Counterfactuals raise a gap concept's mastery and re-evaluate the planner's score, so the
claimed effect of closing a gap is measured rather than asserted. An unmet prerequisite
immediately preceding the recommendation therefore outranks an equal shortfall further
back, and every clause of the generated explanation traces to a number the model produced.

% =====================================================================
\section{Implementation}
\label{sec:impl}

\subsection{Development Tools and Computing Environment}
\label{subsec:tools}

The framework is implemented in Python 3.11 with PyTorch. Data preparation runs in SQL on an
embedded DuckDB database so every transformation is auditable; NetworkX builds the graph and
FP-Growth (\texttt{mlxtend}) mines co-requisites. Training used one NVIDIA T4 GPU and all
inference runs on CPU; data preparation takes about 2.5 minutes and training one
configuration about 25 minutes.

\subsection{Dataset Configuration}
\label{subsec:config}

A pipeline built on SQL, Pandas and NumPy turns the raw records into per-enrolment
sequences, 27{,}904 of them from 25{,}101 learners, with 242{,}670 supervised assessment
positions (Table~\ref{tab:dataset}). Categorical attributes such as gender and education
level are one-hot encoded, and long-tailed click counts are log-transformed before
standardization so a few very active students do not dominate the scale. Sequences are
windowed to the most recent 200 interactions. Cross-validation folds are stratified on final
outcome, and no learner appears in more than one fold.

\subsection{Model Implementation Specifications}
\label{subsec:specs}

The encoder has two Transformer layers, hidden size 128, four heads
and sequence length 200; the three GCN layers use layer normalization and ReLU, giving
64-dimensional concept embeddings. Training uses Adam (learning rate $10^{-3}$, weight decay
$10^{-4}$, cosine annealing), dropout $0.2$, batch size 64 and early stopping at patience 3.
The model has 466{,}569 parameters, and one batched pass over all admissible concepts keeps
an explained recommendation near 160~ms on CPU. Random seeds are fixed and recorded with
every result.

\subsection{Evaluation and Validation}
\label{subsec:eval}

Baselines are a majority-class predictor, SAKT, DKT and a GNN-based recommender. Predictive
quality uses AUC-ROC, RMSE and expected calibration error. Path quality is measured against
observed behaviour: held-out histories are truncated, each planner recommends, and ground
truth is the concepts the learner actually studied next, with only learners who passed or
achieved distinction contributing ground truth. Explanation quality is measured by
fidelity, sufficiency and stability. Ablations remove the knowledge graph and the temporal
decay, under both the conventional protocol and the cold-concept protocol of
Sect.~\ref{sec:results}. Cross-validation is five-fold and grouped by student; comparisons
use paired $t$-tests with Holm--Bonferroni correction \cite{holm1979}, with Cohen's $d$
reported beside every $p$-value, since with many paired decisions significance is attainable
at negligible effect size.

% =====================================================================
\section{Results and Discussion}
\label{sec:results}

\begin{table}[t]
\caption{Predictive performance and component ablation. Folds grouped by student, $p$ Holm-corrected. Proposed expected calibration error $=0.0054$}
\label{tab:accuracy}
\begin{tabular}{p{4.4cm}p{3.1cm}p{1.8cm}p{1.6cm}}
\hline\noalign{\smallskip}
\textbf{Model} & \textbf{AUC-ROC} & \textbf{RMSE} & \textbf{$p$} \\
\noalign{\smallskip}\svhline\noalign{\smallskip}
Majority (no-skill) & $0.5000 \pm 0.0000$ & $0.4669$ & $<0.001$ \\
SAKT                & $0.9233 \pm 0.0018$ & $0.2992$ & $<0.001$ \\
DKT                 & $0.9435 \pm 0.0017$ & $0.2751$ & $0.0009$ \\
GNN-based           & $0.9502 \pm 0.0014$ & $0.2665$ & $0.0220$ \\
\textbf{Proposed}   & $\mathbf{0.9538 \pm 0.0019}$ & $\mathbf{0.2621}$ & --- \\
\noalign{\smallskip}\hline\noalign{\smallskip}
\quad w/o knowledge graph    & $0.9539 \pm 0.0021$ & $0.2618$ & $0.7704$ \\
\quad w/o temporal attention & $0.9532 \pm 0.0014$ & $0.2624$ & $0.4507$ \\
\noalign{\smallskip}\hline\noalign{\smallskip}
\end{tabular}
\end{table}

\begin{table}[t]
\caption{Learning path quality over 906 decisions from 400 held-out learners. Effect size against the proposed planner: syllabus order $d=0.36$, random $d=0.09$; all $p<0.02$}
\label{tab:pathquality}
\begin{tabular}{p{3.6cm}p{2.0cm}p{1.8cm}p{1.8cm}p{1.8cm}}
\hline\noalign{\smallskip}
\textbf{Planner} & \textbf{NDCG@5} & \textbf{Hit@5} & \textbf{Viol.} & \textbf{Cov.} \\
\noalign{\smallskip}\svhline\noalign{\smallskip}
Syllabus order      & $0.0612$ & $0.1832$ & $0.0000$ & $0.452$ \\
Popularity          & $0.1069$ & $0.2936$ & $0.0000$ & $0.510$ \\
Random (admissible) & $0.1300$ & $0.3455$ & $0.0000$ & $0.722$ \\
Weakest-first       & $0.1344$ & $0.3642$ & $0.0000$ & $0.515$ \\
\textbf{Proposed}   & $\mathbf{0.1532}$ & $\mathbf{0.3687}$ & $\mathbf{0.0000}$ & $\mathbf{0.726}$ \\
\noalign{\smallskip}\hline\noalign{\smallskip}
\end{tabular}
\end{table}

\begin{table}[t]
\caption{Cold-concept ablation. Assessments withheld for 27 of the 89 assessed concepts and scored on those only}
\label{tab:coldconcept}
\begin{tabular}{p{4.4cm}p{3.6cm}p{3.0cm}}
\hline\noalign{\smallskip}
\textbf{Variant} & \textbf{AUC-ROC} & \textbf{RMSE} \\
\noalign{\smallskip}\svhline\noalign{\smallskip}
With knowledge graph & $\mathbf{0.9004 \pm 0.0059}$ & $\mathbf{0.3359}$ \\
Without              & $0.8679 \pm 0.0042$ & $0.3611$ \\
\noalign{\smallskip}\hline\noalign{\smallskip}
\textbf{Difference}  & \multicolumn{2}{l}{$\mathbf{+0.0325}$, $p = 0.0003$, $d = 5.06$} \\
\noalign{\smallskip}\hline\noalign{\smallskip}
\end{tabular}
\end{table}

\figslot{fig5.png}{Calibration of predicted mastery across the five folds. Expected
calibration error falls from 0.0084 to 0.0054 after temperature scaling}{fig:calibration}

\subsection{Predictive Accuracy}
\label{subsec:accuracy}

As Table~\ref{tab:accuracy} shows, the proposed model attains the highest AUC-ROC
($0.9538 \pm 0.0019$) and lowest RMSE ($0.2621$), exceeding DKT by $0.0103$
($p = 0.0009$) and the GNN baseline by $0.0036$ ($p = 0.022$). The majority baseline scores
exactly $0.5000$, so the result is not an artefact of class imbalance, and SAKT reaches only
$0.9233$, so the gain is not simply from attention. Calibration error falls from $0.0084$
to $0.0054$ after temperature scaling (Fig.~\ref{fig:calibration}). Because predicted
mastery is shown to learners inside explanations, a poorly calibrated probability would be
an incorrect statement rather than a cosmetic defect.

\subsection{Learning Path Quality}
\label{subsec:pathquality}

Across 906 decisions from 400 held-out successful learners (Table~\ref{tab:pathquality}),
the proposed planner achieves the highest NDCG@5 ($0.1532$) and Hit@5 ($0.3687$). All
differences are significant, but effect sizes qualify this: $d = 0.36$ against syllabus
order and only $d = 0.09$ against random admissible selection. The prerequisite constraint
does most of the work, with the learned ranking adding a small consistent gain. Violations
are zero for every planner, and coverage of $0.726$ shows the planner does not collapse
onto popular concepts, the failure mode popularity-based ranking shows at $0.510$. The
small margin over random selection also bounds the headroom for a learned policy.

\subsection{Explanation Quality}
\label{subsec:explquality}

Explanations were evaluated objectively rather than by participant rating. The sample is
small, 25 decisions with three perturbations each, but every measure is computed against
the model's own counterfactual response rather than an opinion. Fidelity is $0.5318$ and
stability $0.5310$: cited gaps are positively but imperfectly predictive of measured
improvement, and about half persist under perturbation. Sufficiency of $0.0020$ is a
negative finding. The named concept accounted for only $0.2\%$ of predicted benefit, the
rest arising from unmentioned related concepts, so explanations were accurate but
incomplete, and the renderer was amended to name the main beneficiaries. No rating scale
would reveal this, because participants cannot judge a quantity the explanation omits.

\subsection{Why Conventional Ablation Is Insufficient}
\label{subsec:ablation}

Removing the knowledge graph changes accuracy by $-0.0001$ ($p = 0.77$) and removing
temporal attention by $-0.0006$ ($p = 0.45$) (Table~\ref{tab:accuracy}). Read at face
value, neither contributes. But accuracy is computed only at assessment positions, covering
89 concepts with about 2{,}700 labelled examples each, where a free embedding learns each
concept adequately without any structural information, so the graph is genuinely redundant
there. The graph exists for the other 148 concepts, which never enter the metric.

To measure it directly, all assessments were withheld for 27 of the 89 assessed concepts
and their responses removed from the input; evaluation was restricted to those concepts.
Table~\ref{tab:coldconcept} shows $0.9004$ with the graph against $0.8679$ without, a gain
of $+0.0325$ ($p = 0.0003$, $d = 5.06$) on every fold. The graph contributes nothing where
supervision is plentiful and substantially where it is absent, which matters because the
recommender ranks across all 237 concepts. Component ablations should be accompanied by
evidence that the metric can detect the component removed.

\subsection{Overall Discussion}
\label{subsec:discussion}

Combining knowledge graph modeling, temporal knowledge tracing, constrained planning and
explanation improves both predictive and recommendation metrics over the baselines.
Explanation adds computation, about 160~ms for a fully explained recommendation against
7~ms for a cached mastery lookup, which remains well within the latency of an interactive
advising interface.

% =====================================================================
\section{Conclusion and Future Work}
\label{sec:conclusion}

\subsection{Conclusion}
\label{subsec:conclusion}

This study combined knowledge graph modeling, Deep Knowledge Tracing,
prerequisite-constrained planning and priority-backtracking explanation into an explainable
learning path recommender. On OULAD it exceeded majority, SAKT, DKT and GNN baselines
(AUC-ROC $0.9538$, calibration error $0.0054$) and made no prerequisite violations. Its
principal finding is methodological: a knowledge graph that conventional ablation reports as
useless improves cold-concept accuracy by $+0.0325$ with $d = 5.06$. The discrepancy is a
property of the metric, not the model: aggregate accuracy measured where supervision is
abundant cannot observe a structural prior whose purpose is to compensate for its absence.
Verifying such systems therefore requires evaluation protocols matched to what each
component is for.

\subsection{Limitations}
\label{subsec:limitations}

The prediction target is successful completion rather than knowledge. Because OULAD records
only submitted assessments, treating non-submission as failure is necessary for a usable
class balance, but a substantial share of the signal is then engagement rather than
understanding: appropriate for a system intended to prompt intervention, but a weaker claim
than measuring mastery. The concept graph is likewise derived rather than given. OULAD
supplies no concept annotations or material titles, placement for $82\%$ of materials
relies on observed access, and mining prerequisites from the same behaviour that orders
the concepts is partially self-referential.

Supervision covers 89 of 237 concepts, so mastery for the remainder is not directly
calibrated, and predicted mastery is bimodal, which is why learner-relative thresholds are
needed. The proposed model uses learner features that the DKT and SAKT baselines, following
their original formulations, do not, so part of the gain may reflect feature access rather
than architecture. Finally, all data comes from a single institution's distance-learning
courses, and generalization to campus-based teaching is untested.

\subsection{Ethical Considerations}
\label{subsec:ethics}

OULAD is publicly released under CC-BY 4.0 and anonymised, and no data was collected from
participants. The system is advisory. Because the model conditions on gender, disability,
deprivation band and age band, it should not inform assessment or progression decisions
until a subgroup fairness analysis has been carried out.

\subsection{Future Work}
\label{subsec:future}

Four directions follow from the limitations. The most immediate is a \textbf{fairness
analysis}: subgroup performance across gender, disability, deprivation band and age band
must be established before any deployment. Second, a \textbf{learned policy} over the same
state and reward should be trained and evaluated, in simulation and by off-policy
estimation on logged trajectories; algorithms such as Soft Actor-Critic or Proximal Policy
Optimization would need to beat a greedy planner already $d = 0.09$ above random selection
to justify their complexity. Third, a \textbf{user evaluation} would address what the
objective measures cannot: fidelity, sufficiency and stability show whether explanations
are faithful to the model, not whether learners find them useful. Fourth, \textbf{expert
validation of the mined edges} would give an independent estimate of graph precision.
Beyond these, integrating forum and video-lecture activity, testing transfer across
subjects and institutions, and building lighter models through pruning and distillation
are natural extensions.

\begin{acknowledgement}
All source code, SQL transformations and results are publicly available, and every table in
this chapter is regenerated from the raw data by a single command with fixed random seeds:
\url{https://github.com/abhay27singh/explainable-learning-path-recommender}.
\end{acknowledgement}

% =====================================================================
\begin{thebibliography}{99}

\bibitem{li2026survey}
Li, A., Li, Y., Gao, X.: Personalized learning path recommendation based on knowledge
graphs: a survey. Electronics \textbf{15}(1), 238 (2026)

\bibitem{aher2013combination}
Aher, S.B., Lobo, L.M.R.J.: Combination of machine learning algorithms for recommendation
of courses in e-learning system based on historical data. Knowl.-Based Syst. \textbf{51},
1--14 (2013)

\bibitem{kuzilek2017oulad}
Kuzilek, J., Hlosta, M., Zdrahal, Z.: Open University Learning Analytics dataset. Sci. Data
\textbf{4}, 170171 (2017)

\bibitem{corbett1994knowledge}
Corbett, A.T., Anderson, J.R.: Knowledge tracing: modeling the acquisition of procedural
knowledge. User Model. User-Adapt. Interact. \textbf{4}(4), 253--278 (1994)

\bibitem{piech2015dkt}
Piech, C., Bassen, J., Huang, J., Ganguli, S., Sahami, M., Guibas, L., Sohl-Dickstein, J.:
Deep knowledge tracing. In: Advances in Neural Information Processing Systems, vol.~28,
pp.~505--513 (2015)

\bibitem{zhang2017dkvmn}
Zhang, J., Shi, X., King, I., Yeung, D.-Y.: Dynamic key-value memory networks for knowledge
tracing. In: Proceedings of the 26th International Conference on World Wide Web,
pp.~765--774 (2017)

\bibitem{ghosh2020akt}
Ghosh, A., Heffernan, N.T., Lan, A.S.: Context-aware attentive knowledge tracing. In:
Proceedings of the 26th ACM SIGKDD International Conference on Knowledge Discovery and Data
Mining, pp.~2330--2339 (2020)

\bibitem{shen2021lpkt}
Shen, S., Liu, Q., Chen, E., Huang, Z., Huang, W., Yin, Y., Su, Y., Wang, S.: Learning
process-consistent knowledge tracing. In: Proceedings of the 27th ACM SIGKDD Conference on
Knowledge Discovery and Data Mining, pp.~1452--1460 (2021)

\bibitem{vaswani2017attention}
Vaswani, A., Shazeer, N., Parmar, N., Uszkoreit, J., Jones, L., Gomez, A.N., Kaiser, L.,
Polosukhin, I.: Attention is all you need. In: Advances in Neural Information Processing
Systems, vol.~30, pp.~5998--6008 (2017)

\bibitem{kipf2017gcn}
Kipf, T.N., Welling, M.: Semi-supervised classification with graph convolutional networks.
In: International Conference on Learning Representations (2017)

\bibitem{mnih2015dqn}
Mnih, V., Kavukcuoglu, K., Silver, D., Rusu, A.A., Veness, J., Bellemare, M.G., Graves, A.,
Riedmiller, M., Fidjeland, A.K., Ostrovski, G., et al.: Human-level control through deep
reinforcement learning. Nature \textbf{518}, 529--533 (2015)

\bibitem{khosravi2022xaied}
Khosravi, H., Buckingham Shum, S., Chen, G., Conati, C., Tsai, Y.-S., Kay, J., Knight, S.,
Martinez-Maldonado, R., Sadiq, S., Ga\v{s}evi\'{c}, D.: Explainable artificial intelligence
in education. Comput. Educ.: Artif. Intell. \textbf{3}, 100074 (2022)

\bibitem{ribeiro2016lime}
Ribeiro, M.T., Singh, S., Guestrin, C.: ``Why should I trust you?'' Explaining the
predictions of any classifier. In: Proceedings of the 22nd ACM SIGKDD International
Conference on Knowledge Discovery and Data Mining, pp.~1135--1144 (2016)

\bibitem{lundberg2017shap}
Lundberg, S.M., Lee, S.-I.: A unified approach to interpreting model predictions. In:
Advances in Neural Information Processing Systems, vol.~30, pp.~4765--4774 (2017)

\bibitem{zhang2023finegrained}
Zhang, S., Liu, Q., Chen, Y., Wang, L., Zhao, H.: A fine-grained and multi-context-aware
learning path recommendation model over knowledge graphs for online learning communities.
Inf. Process. Manag. \textbf{60}(5), 103464 (2023)

\bibitem{yang2025explainable}
Yang, Y., Peng, X., Chen, M., Liu, S.: An explainable graph-based course recommendation
model based on multiple interest factors. Expert Syst. Appl. \textbf{264}, 125889 (2025)

\bibitem{xie2025profile}
Xie, X., Feng, X.: Personalized learning path recommendation based on learner profile and
knowledge graph. Comput. Inform. \textbf{44}(4), 983--1008 (2025)

\bibitem{zhao2026rlkt}
Zhao, J., Wang, H., Li, M., Sun, Y.: Personalized learning path recommendation for middle
school students based on joint optimization of deep reinforcement learning and knowledge
tracing. Sci. Rep. (2026)

\bibitem{boujmiraz2026review}
Boujmiraz, S., Darhmaoui, H., Drissi El Maliani, A.: Predicting student performance: a
comprehensive review of machine learning, deep learning, and explainable AI approaches.
Comput. Educ.: Artif. Intell. \textbf{10}, 100548 (2026)

\bibitem{qiang2026xaiedu}
Qiang, M., Liu, Z., Zhang, R.: Explainable AI in education: integrating educational domain
knowledge into the deep learning model for improved student performance prediction. Sci.
Rep. \textbf{16}, 9515 (2026)

\bibitem{pandey2019sakt}
Pandey, S., Karypis, G.: A self-attentive model for knowledge tracing. In: Proceedings of
the 12th International Conference on Educational Data Mining, pp.~384--389 (2019)

\bibitem{guo2017calibration}
Guo, C., Pleiss, G., Sun, Y., Weinberger, K.Q.: On calibration of modern neural networks.
In: Proceedings of the 34th International Conference on Machine Learning, pp.~1321--1330
(2017)

\bibitem{holm1979}
Holm, S.: A simple sequentially rejective multiple test procedure. Scand. J. Stat.
\textbf{6}(2), 65--70 (1979)

\end{thebibliography}

\end{document}
